%% ========================================================================
%% Chapter 3: Dasar Input/Output (I/O)
%% ========================================================================

\chapter{Dasar Input/Output (I/O)}
\label{ch:basic-io}

\chapterhero{violet}{Pelajari bagaimana aliran data bekerja di shell melalui file descriptor, redirection, pipeline, dan teknik pemrosesan input-output yang efisien.}{\faExchangeAlt\quad stdin/stdout}{\faProjectDiagram\quad pipeline}{\faSignInAlt\quad redirection}

Bayangkan kamu bekerja di warung fotokopi. Ada tiga jalur kerja: jalur masuk untuk menerima dokumen dari pelanggan, jalur keluar untuk menyerahkan hasil fotokopi, dan jalur khusus untuk laporan kesalahan jika mesin rusak. Linux menggunakan prinsip yang sama untuk setiap program: ada saluran khusus untuk input, output biasa, dan pesan kesalahan. Inilah konsep dasar I/O di Linux.

Filosofi Unix menyatakan bahwa setiap program hanya melakukan satu hal dengan baik, dan program-program kecil itu bisa digabungkan untuk menyelesaikan tugas yang kompleks. Penggabungan ini dimungkinkan karena semua program berbagi antarmuka I/O yang seragam. Output dari satu program bisa langsung menjadi input program berikutnya, seperti jalur produksi di pabrik.

Pemahaman tentang I/O adalah syarat untuk hampir semua pekerjaan administrasi Linux. Menyimpan log, memfilter output, memisahkan pesan kesalahan dari data normal, semua bergantung pada mekanisme yang dipelajari di bab ini.

\textbf{Yang Akan Kamu Pelajari}

Setelah menyelesaikan bab ini, kamu akan mampu:

\begin{enumerate}
    \item menjelaskan konsep file descriptor dan mengidentifikasi stdin, stdout, serta stderr;
    \item mengalihkan output ke berkas menggunakan operator \texttt{>} dan \texttt{>>} serta memisahkan stderr dengan \texttt{2>};
    \item membangun pipeline bertingkat menggunakan operator \texttt{|} dengan kombinasi \cmd{grep}, \cmd{sort}, \cmd{uniq}, dan \cmd{wc};
    \item menggunakan \cmd{tee} untuk mengirim output ke terminal dan berkas secara bersamaan;
    \item menerapkan here-document untuk membuat berkas konten multi-baris langsung dari baris perintah.
\end{enumerate}

\begin{notebox}
Seluruh praktek pada bab ini dijalankan pada direktori \dir{~/lab-os/chapter03-io/}.
Buat direktori ini sebelum memulai praktek pertama:
\begin{lstlisting}[language=bash]
mkdir -p ~/lab-os/chapter03-io
cd ~/lab-os/chapter03-io
\end{lstlisting}
\end{notebox}

%% ------------------------------------------------------------------------
\section{File Descriptor}
\label{sec:file-descriptors}

Setiap kali program membuka berkas, kernel memberikan angka bulat yang disebut \textit{file descriptor} (FD) sebagai pegangan untuk operasi baca-tulis. Angka ini ringan dan efisien: program cukup menyebut angkanya, dan kernel tahu berkas mana yang dimaksud.

Tiga file descriptor pertama sudah tersedia otomatis untuk setiap proses saat pertama kali dijalankan. File descriptor 0 adalah \textit{stdin} (standard input), yaitu saluran untuk membaca input. Secara default stdin terhubung ke keyboard. File descriptor 1 adalah \textit{stdout} (standard output), yaitu saluran untuk menulis output normal. Secara default stdout terhubung ke layar terminal. File descriptor 2 adalah \textit{stderr} (standard error), yaitu saluran khusus untuk pesan kesalahan. Stderr juga terhubung ke layar secara default, tapi bisa dialihkan secara terpisah dari stdout.

Pemisahan antara stdout dan stderr bukan sekadar konvensi. Ketika kamu membangun pipeline seperti \texttt{perintah1 | perintah2}, hanya stdout dari \texttt{perintah1} yang mengalir ke stdin \texttt{perintah2}. Pesan kesalahan dari \texttt{perintah1} tidak ikut mengalir sehingga tidak mencemari data yang diproses \texttt{perintah2}.

Linux menyimpan informasi tentang file descriptor yang terbuka oleh setiap proses di direktori \dir{/proc/PID/fd}. Di sini, variabel shell \texttt{\$\$} merujuk ke PID proses shell yang sedang berjalan.

\subsection*{Praktek 3.1: Amati File Descriptor Proses Aktif}

Tujuan praktek ini adalah melihat file descriptor yang dimiliki proses shell dan memahami kemana masing-masing mengarah.

\begin{enumerate}
    \item Masuk ke direktori kerja:
\begin{lstlisting}[language=bash, caption={Masuk ke direktori kerja}]
cd ~/lab-os/chapter03-io
\end{lstlisting}

    \item Lihat file descriptor yang dimiliki shell saat ini:
\begin{lstlisting}[language=bash, caption={File descriptor proses shell aktif}]
ls -l /proc/$$/fd
\end{lstlisting}
    \textit{Amati:} Tiga entri pertama (0, 1, 2) biasanya mengarah ke \file{/dev/pts/0} atau nama serupa. Itu adalah pseudo-terminal yang mewakili sesi terminal kamu. Ketiganya mengarah ke tempat yang sama karena input, output, dan error semuanya terhubung ke layar yang sama.

    \item Lihat PID shell saat ini:
\begin{lstlisting}[language=bash, caption={PID shell saat ini}]
echo "PID shell ini: $$"
\end{lstlisting}

    \item Buka berkas untuk dibaca dan periksa file descriptor baru yang muncul:
\begin{lstlisting}[language=bash, caption={Amati penambahan file descriptor baru}]
exec 3< /etc/hostname
ls -l /proc/$$/fd
exec 3<&-
\end{lstlisting}
    \textit{Amati:} Setelah menjalankan \texttt{exec 3<}, file descriptor 3 muncul di daftar mengarah ke \file{/etc/hostname}. Setelah \texttt{exec 3<\&-}, file descriptor itu ditutup dan hilang dari daftar.

    \item Cek file descriptor proses lain (misalnya init):
\begin{lstlisting}[language=bash, caption={File descriptor proses PID 1}]
sudo ls -l /proc/1/fd | head -10
\end{lstlisting}
\end{enumerate}

\begin{challengebox}
Alihkan hanya stderr ke berkas, sementara stdout tetap tampil di terminal. Jalankan perintah \cmd{find /etc -name "*.conf"} yang akan menghasilkan dua jenis output: daftar berkas yang ditemukan (stdout) dan pesan "Permission denied" (stderr). Pisahkan keduanya dengan: \texttt{find /etc -name "*.conf" 2> error-find.txt}. Periksa isi \file{error-find.txt} dengan \cmd{cat error-find.txt}. Berapa banyak kesalahan izin akses yang tercatat?
\end{challengebox}

%% ------------------------------------------------------------------------
\section{Pengalihan Output}
\label{sec:io-redirection}

Operator pengalihan output memberi instruksi ke shell untuk menghubungkan stdout atau stderr ke berkas alih-alih ke terminal. Ini adalah cara paling mendasar untuk menyimpan hasil perintah.

Operator \texttt{>} mengalihkan stdout ke berkas. Jika berkas sudah ada, isinya ditimpa seluruhnya. Jika belum ada, berkas baru dibuat. Operator \texttt{>>} juga mengalihkan stdout ke berkas, tapi dalam mode tambah: isi lama dipertahankan dan output baru ditambahkan di akhir. Perbedaan ini krusial: salah menggunakan \texttt{>} pada berkas log bisa menghapus riwayat yang tidak bisa dipulihkan.

Untuk stderr, gunakan \texttt{2>} (overwrite) atau \texttt{2>>} (tambah). Angka 2 adalah nomor file descriptor untuk stderr. Jika ingin menggabungkan stderr ke aliran yang sama dengan stdout, gunakan \texttt{2>\&1}: ini berarti "alihkan FD 2 ke tempat yang sama dengan FD 1". Urutan sangat penting: tulis \texttt{> berkas 2>\&1}, bukan sebaliknya.

Berkas khusus \file{/dev/null} berfungsi sebagai tempat sampah. Semua yang ditulis ke sana langsung lenyap. Gunakan ini untuk membuang output yang tidak perlu, seperti pesan "Permission denied" yang tidak relevan saat mencari berkas di seluruh sistem.

\subsection*{Praktek 3.2: Simpan Output Perintah ke Berkas}

Tujuan praktek ini adalah berlatih menggunakan operator \texttt{>}, \texttt{>>}, \texttt{2>}, dan \texttt{2>\&1} dalam situasi nyata.

\begin{enumerate}
    \item Simpan daftar berkas di direktori kerja ke berkas:
\begin{lstlisting}[language=bash, caption={Simpan output ke berkas baru}]
ls -lh ~/lab-os/ > daftar-lab.txt
cat daftar-lab.txt
\end{lstlisting}
    \textit{Amati:} Output tidak tampil di terminal karena sudah dialihkan ke berkas. Berkas \file{daftar-lab.txt} berisi hasil yang sama persis dengan yang biasanya tampil di layar.

    \item Tambahkan informasi sistem ke berkas yang sama tanpa menghapus isi lama:
\begin{lstlisting}[language=bash, caption={Tambahkan output ke berkas yang sudah ada}]
echo "--- Info Sistem ---" >> daftar-lab.txt
uname -a >> daftar-lab.txt
date >> daftar-lab.txt
cat daftar-lab.txt
\end{lstlisting}

    \item Buat berkas log yang memisahkan stdout dan stderr:
\begin{lstlisting}[language=bash, caption={Pisahkan stdout dan stderr ke berkas berbeda}]
find /etc -name "*.conf" > konfig-ditemukan.txt 2> konfig-error.txt
echo "Berkas ditemukan: $(wc -l < konfig-ditemukan.txt)"
echo "Baris error: $(wc -l < konfig-error.txt)"
\end{lstlisting}
    \textit{Amati:} Dua berkas terpisah terbentuk: satu berisi path berkas yang berhasil ditemukan, satu lagi berisi pesan kesalahan izin akses.

    \item Buang semua pesan error dan simpan hanya output yang berhasil:
\begin{lstlisting}[language=bash, caption={Buang stderr ke /dev/null}]
find /etc -name "*.conf" 2>/dev/null > konfig-bersih.txt
wc -l konfig-bersih.txt
\end{lstlisting}

    \item Gabungkan stdout dan stderr ke satu berkas log:
\begin{lstlisting}[language=bash, caption={Gabungkan stdout dan stderr ke satu berkas}]
find /var/log -name "*.log" > semua-output.txt 2>&1
head -20 semua-output.txt
\end{lstlisting}
    \textit{Amati:} Berkas \file{semua-output.txt} berisi campuran path berkas yang ditemukan dan pesan error, sesuai urutan kemunculannya.
\end{enumerate}

\begin{challengebox}
Pisahkan stdout dan stderr ke dua berkas berbeda dengan satu perintah. Jalankan \cmd{ls /etc /direktori-tidak-ada} (berikan satu direktori yang ada dan satu yang tidak ada). Gunakan: \texttt{ls /etc /direktori-tidak-ada > hasil.txt 2> error.txt}. Periksa isi keduanya. Kemudian ulangi dengan menggabungkan keduanya menggunakan \texttt{\&>}: \texttt{ls /etc /direktori-tidak-ada \&> semua.txt}. Bandingkan isi \file{semua.txt} dengan gabungan manual \file{hasil.txt} dan \file{error.txt}.
\end{challengebox}

%% ------------------------------------------------------------------------
\section{Pipe dan Pipeline}
\label{sec:pipes}

\term{Pipe} adalah mekanisme yang menghubungkan stdout satu perintah langsung ke stdin perintah berikutnya, tanpa perlu berkas perantara. Operator pipe adalah \texttt{|}. Proses dalam pipeline berjalan bersamaan: perintah di sebelah kiri mulai menghasilkan output, dan perintah di sebelah kanan langsung memprosesnya tanpa menunggu perintah kiri selesai sepenuhnya.

Kekuatan \term{pipe} terletak pada komposisi. Setiap perintah Unix melakukan satu tugas dengan baik. \cmd{sort} mengurutkan baris. \cmd{uniq} menghapus baris duplikat berurutan (flag \texttt{-c} menambahkan hitungan). \cmd{wc} menghitung baris, kata, atau karakter. \cmd{head} dan \cmd{tail} mengambil bagian tertentu dari output. Kombinasi perintah-perintah sederhana ini bisa menjawab pertanyaan yang kompleks.

Seperti yang sudah dipelajari di Bab~2 dengan \cmd{grep} dan \cmd{awk}, kemampuan pipeline semakin berguna ketika digabungkan dengan perintah manipulasi teks. Misalnya, untuk menemukan sepuluh IP address yang paling sering muncul di log web: ambil semua baris dengan \cmd{cat}, ekstrak kolom IP dengan \cmd{awk}, urutkan dengan \cmd{sort}, hitung kemunculan dengan \cmd{uniq -c}, urutkan lagi secara numerik dengan \cmd{sort -rn}, dan ambil sepuluh teratas dengan \cmd{head -10}. Semua dalam satu baris perintah.

Pipeline panjang bisa dibagi ke beberapa baris menggunakan backslash \texttt{\textbackslash} di akhir baris untuk keterbacaan.

\subsection*{Praktek 3.3: Bangun Pipeline Bertingkat}

Tujuan praktek ini adalah membangun pipeline langkah demi langkah untuk menganalisis data sistem.

\begin{enumerate}
    \item Mulai dengan perintah sederhana, lalu tambahkan tahap satu per satu:
\begin{lstlisting}[language=bash, caption={Pipeline bertahap: hitung pengguna per shell}]
cat /etc/passwd
\end{lstlisting}
    \textit{Amati:} Setiap baris adalah satu entri pengguna dengan kolom dipisahkan titik dua.

    \item Ekstrak kolom terakhir (shell yang digunakan) menggunakan \cmd{awk}:
\begin{lstlisting}[language=bash, caption={Ekstrak kolom shell dari /etc/passwd}]
cat /etc/passwd | awk -F: '{print $NF}'
\end{lstlisting}

    \item Urutkan hasilnya:
\begin{lstlisting}[language=bash, caption={Urutkan daftar shell}]
cat /etc/passwd | awk -F: '{print $NF}' | sort
\end{lstlisting}

    \item Hitung berapa kali setiap shell muncul dan urutkan dari terbanyak:
\begin{lstlisting}[language=bash, caption={Hitung kemunculan setiap shell}]
cat /etc/passwd | awk -F: '{print $NF}' | sort | uniq -c | sort -rn
\end{lstlisting}
    \textit{Amati:} Output menampilkan hitungan di sebelah kiri dan nama shell di sebelah kanan, diurutkan dari yang paling banyak digunakan.

    \item Analisis penggunaan disk: tampilkan lima direktori terbesar di \dir{/var}:
\begin{lstlisting}[language=bash, caption={Lima direktori terbesar di /var}]
sudo du -sh /var/* 2>/dev/null | sort -h | tail -5
\end{lstlisting}
    \textit{Amati:} Flag \texttt{-h} pada \cmd{sort} menangani satuan seperti K, M, G dengan benar. Tanpa flag ini, \texttt{10M} bisa dianggap lebih kecil dari \texttt{9G}.

    \item Hitung jumlah proses yang berjalan per pengguna:
\begin{lstlisting}[language=bash, caption={Jumlah proses per pengguna}]
ps aux | tail -n +2 | awk '{print $1}' | sort | uniq -c | sort -rn
\end{lstlisting}
    \textit{Amati:} \texttt{tail -n +2} melewatkan baris header \cmd{ps aux}. Hasilnya menunjukkan pengguna mana yang paling banyak menjalankan proses.
\end{enumerate}

\begin{challengebox}
Bangun pipeline yang menggabungkan teknik dari Bab~2 (grep dan awk) dengan sort dan uniq. Analisis berkas \file{/var/log/auth.log} (atau \file{/var/log/syslog} jika auth.log tidak tersedia). Cari baris yang mengandung kata \texttt{Failed} menggunakan \cmd{grep}, ekstrak kata kunci di kolom tertentu menggunakan \cmd{awk}, hitung kemunculan dengan \cmd{sort | uniq -c}, lalu urutkan hasilnya. Jika kedua berkas tidak ada, buat berkas latihan sendiri dengan beberapa baris yang mengandung kata Failed dan analisis menggunakan pipeline yang sama.
\end{challengebox}

%% ------------------------------------------------------------------------
\section{tee dan Teknik Lanjutan}
\label{sec:tee-advanced}

Ketika kamu mengalihkan output ke berkas dengan \texttt{>}, output tidak lagi tampil di terminal. Ini bermasalah saat menjalankan perintah panjang yang ingin dipantau secara langsung sekaligus disimpan. Perintah \cmd{tee} memecahkan masalah ini: ia membaca dari stdin dan menulis ke stdout sekaligus ke berkas. Namanya diambil dari bentuk huruf T yang menggambarkan aliran data bercabang dua.

Penggunaan dasarnya: \texttt{perintah | tee berkas.log}. Untuk mode tambah agar berkas lama tidak tertimpa, tambahkan flag \texttt{-a}: \texttt{perintah | tee -a berkas.log}. Kamu juga bisa menulis ke beberapa berkas sekaligus: \texttt{perintah | tee berkas1.txt berkas2.txt}.

Here-document (\texttt{<<EOF}) adalah cara memberi input multi-baris ke perintah langsung dari baris perintah, tanpa membuat berkas terpisah terlebih dahulu. Shell membaca semua baris antara penanda pembuka (\texttt{<<EOF}) dan penanda penutup (\texttt{EOF}) sebagai input. Penanda bisa berupa kata apa saja, asalkan pembuka dan penutup sama. Jika penanda pembuka diapit tanda kutip tunggal (\texttt{'EOF'}), variabel di dalam blok tidak akan diekspansi.

Here-document sangat berguna untuk membuat berkas konfigurasi dalam skrip, menulis pesan multi-baris, atau memberikan input ke perintah seperti \cmd{cat} dan \cmd{tee}.

\subsection*{Praktek 3.4: Output ke Berkas dan Terminal Sekaligus}

Tujuan praktek ini adalah menggunakan \cmd{tee} dalam pipeline dan membuat berkas dengan here-document.

\begin{enumerate}
    \item Gunakan \cmd{tee} untuk melihat sekaligus menyimpan output:
\begin{lstlisting}[language=bash, caption={tee: output ke terminal dan berkas}]
df -h | tee laporan-disk.txt
echo "---"
cat laporan-disk.txt
\end{lstlisting}
    \textit{Amati:} Output \cmd{df -h} tampil di terminal saat perintah dijalankan. Berkas \file{laporan-disk.txt} menyimpan konten yang sama persis.

    \item Gunakan \cmd{tee} dalam pipeline panjang untuk "menyadap" data di tengah:
\begin{lstlisting}[language=bash, caption={tee di tengah pipeline}]
ps aux | tee semua-proses.txt | grep -i ssh
\end{lstlisting}
    \textit{Amati:} Berkas \file{semua-proses.txt} menyimpan semua proses (output penuh dari \cmd{ps aux}), tetapi terminal hanya menampilkan baris yang mengandung ssh.

    \item Tambahkan entri log dengan timestamp menggunakan \cmd{tee -a}:
\begin{lstlisting}[language=bash, caption={Log berulang dengan tee -a}]
echo "$(date '+%Y-%m-%d %H:%M:%S') - pemeriksaan disk" | tee -a aktivitas.log
df -h | grep -v tmpfs | tee -a aktivitas.log
echo "$(date '+%Y-%m-%d %H:%M:%S') - selesai" | tee -a aktivitas.log
cat aktivitas.log
\end{lstlisting}

    \item Buat berkas konfigurasi menggunakan here-document:
\begin{lstlisting}[language=bash, caption={Buat berkas dengan here-document}]
cat << 'EOF' > aplikasi.conf
# Konfigurasi aplikasi
HOST=localhost
PORT=8080
MODE=production
LOG_LEVEL=info
EOF
cat aplikasi.conf
\end{lstlisting}
    \textit{Amati:} Tanda kutip tunggal di \texttt{'EOF'} mencegah variabel seperti \texttt{\$HOST} diekspansi. Coba hapus tanda kutip dan lihat perbedaannya.

    \item Gunakan here-document dengan \cmd{tee} untuk menulis berkas dengan hak superuser:
\begin{lstlisting}[language=bash, caption={Tulis berkas dengan sudo menggunakan tee}]
cat << 'EOF' | sudo tee /tmp/test-sudo.conf > /dev/null
# Berkas uji
tes=berhasil
EOF
cat /tmp/test-sudo.conf
\end{lstlisting}
    \textit{Amati:} Operator \texttt{>} diproses oleh shell sebelum \cmd{sudo}, sehingga tidak mendapat hak akses root. Kombinasi \cmd{tee} dengan \cmd{sudo} adalah cara yang benar untuk menulis berkas yang memerlukan hak superuser.

    \item Buat log berulang bertanggal menggunakan pipeline lengkap:
\begin{lstlisting}[language=bash, caption={Log bertanggal dengan pipeline dan tee}]
LOGFILE="laporan-$(date +%Y%m%d).log"
{
    echo "=== Laporan: $(date) ==="
    echo "--- Disk ---"
    df -h | grep -v tmpfs
    echo "--- Memori ---"
    free -h
    echo "--- Proses Teratas ---"
    ps aux | sort -k3 -rn | head -5
} | tee "$LOGFILE"
echo "Log disimpan di: $LOGFILE"
\end{lstlisting}
\end{enumerate}

\begin{challengebox}
Bangun pipeline lengkap yang menggabungkan teknik dari Bab~2 dan bab ini. Tujuan: buat laporan proses lengkap yang tersimpan dan tampil di terminal sekaligus. Gunakan \cmd{ps aux} untuk mengambil semua proses, saring dengan \cmd{grep} untuk mencari proses milik pengguna kamu saja (\cmd{grep \$USER}), urutkan berdasarkan penggunaan CPU dengan \cmd{sort -k3 -rn}, ambil 10 teratas dengan \cmd{head -10}, dan simpan ke berkas \file{top-proses.log} menggunakan \cmd{tee -a} agar log tidak tertimpa jika dijalankan beberapa kali. Tambahkan timestamp di awal dengan \cmd{echo "\$(date)"} sebelum pipeline utama, juga dialihkan ke berkas yang sama.
\end{challengebox}

%% ------------------------------------------------------------------------
\section{Rangkuman}

Dalam bab ini, kamu telah mempelajari:

\begin{itemize}[noitemsep,topsep=2pt]
    \item \textbf{File descriptor}: setiap proses mendapat tiga FD otomatis: stdin (0), stdout (1), dan stderr (2). Ketiganya bisa dialihkan secara terpisah.
    \item \textbf{Operator \texttt{>}}: mengalihkan stdout ke berkas dengan menimpa isi lama. Gunakan \texttt{>>} untuk menambahkan tanpa menghapus konten sebelumnya.
    \item \textbf{Pemisahan stderr}: operator \texttt{2>} mengalihkan hanya pesan kesalahan ke berkas terpisah. \texttt{2>\&1} menggabungkan stderr ke aliran stdout.
    \item \textbf{/dev/null}: berkas khusus yang membuang semua input. Berguna untuk menekan pesan kesalahan yang tidak relevan.
    \item \textbf{Pipeline}: operator \texttt{|} menghubungkan stdout satu perintah ke stdin berikutnya. Proses berjalan bersamaan dan data mengalir tanpa berkas perantara.
    \item \textbf{tee}: mengirim output ke terminal sekaligus menyimpan ke berkas. Flag \texttt{-a} untuk mode tambah. Bisa ditempatkan di mana saja dalam pipeline.
    \item \textbf{Here-document}: blok \texttt{<\,<EOF ... EOF} memberi input multi-baris langsung ke perintah tanpa perlu berkas sementara.
\end{itemize}

%% ------------------------------------------------------------------------
\section{Latihan}

\textbf{Instruksi Umum}: Kerjakan seluruh latihan secara mandiri. Catat langkah penting, simpan tangkapan layar bila diperlukan, lalu rangkum hasilnya sebagai dokumentasi pribadi.

\subsubsection*{Latihan 3.1 Pengalihan dan Pemisahan I/O}
Jalankan perintah \cmd{find /usr -name "*.so" -size +1M} yang akan menghasilkan dua jenis output: path berkas yang ditemukan (stdout) dan pesan "Permission denied" (stderr). Lakukan tiga hal: (a) simpan stdout ke \file{library-besar.txt} dan stderr ke \file{akses-error.txt} dengan satu perintah, (b) hitung jumlah baris masing-masing berkas menggunakan \cmd{wc -l}, (c) ulangi perintah tapi kali ini buang stderr ke \file{/dev/null} dan simpan hanya stdout. Sertakan tangkapan layar dan jelaskan perbedaan hasil.


\subsubsection*{Latihan 3.2 Pipeline Analisis Sistem}
Bangun pipeline untuk menjawab dua pertanyaan tentang sistem: (a) shell mana yang paling banyak digunakan oleh akun di sistem? Gunakan \file{/etc/passwd} sebagai sumber data, ekstrak kolom terakhir dengan \cmd{awk -F: '\{print \$NF\}'}, lalu hitung dengan \cmd{sort | uniq -c | sort -rn}. (b) Berapa jumlah proses aktif untuk setiap pengguna? Gunakan \cmd{ps aux}, lewatkan baris header, ekstrak nama pengguna, dan hitung kemunculannya. Simpan kedua hasil ke berkas terpisah menggunakan operator \texttt{>}.


\subsubsection*{Latihan 3.3 Membuat Log Sistem Otomatis}
Buat skrip log sederhana langsung dari terminal menggunakan teknik yang dipelajari. Gunakan blok \texttt{\{ ... \}} yang berisi beberapa perintah: timestamp dengan \cmd{date}, penggunaan disk dengan \cmd{df -h}, penggunaan memori dengan \cmd{free -h}, dan lima proses dengan CPU tertinggi menggunakan \cmd{ps aux | sort -k3 -rn | head -5}. Alirkan seluruh blok ke \cmd{tee -a laporan-sistem.log} agar output tampil di terminal sekaligus tersimpan. Jalankan dua kali dan verifikasi bahwa berkas log berisi dua entri berurutan, bukan ditimpa.


%% ------------------------------------------------------------------------
\section{Referensi}

\begin{itemize}
    \item Shotts, W. \textit{The Linux Command Line}, edisi ke-2. No Starch Press, 2019.
    \item Cooper, M. \textit{Advanced Bash-Scripting Guide}. The Linux Documentation Project, 2014.
    \item Ward, B. \textit{How Linux Works}, edisi ke-3. No Starch Press, 2021.
    \item Nemeth, E., Snyder, G., Hein, T. R. \textit{UNIX and Linux System Administration Handbook}, edisi ke-5. Pearson, 2017.
\end{itemize}
